% Journal synthesis paper -- STRUCTURE SKELETON.
% Venue (planned): IJDAR primary; IEEE Access safety; DAS if conference.
% Final camera-ready will use the venue class (Springer svjour3 for IJDAR).
% This skeleton uses 'article' so the structure builds anywhere.
%
% HONESTY RULES BAKED INTO THIS FILE:
%  - \priorfact{...}  : a result ESTABLISHED in the two prior ASYU papers,
%                       cited, NOT re-claimed here. Renders normally.
%  - \pending{...}    : a NEW result this paper will report once the
%                       experiment is run. Renders as a visible PENDING tag.
%                       NEVER replace with a number until the experiment exists.
%  - No long em-dash glyphs; use colons/commas.
% IJDAR / Springer Nature layout EMULATION.
% The genuine Springer class (sn-jnl.cls) is not installed in this
% environment, so this skeleton emulates the IJDAR look with 'article':
% single-column masthead/title/abstract/keywords, then two-column body.
% The real submission swaps in sn-jnl.cls; content is class-independent.
% HONESTY: no real DOI, no acceptance dates, no copyright-of-record line is
% fabricated. The masthead carries only public journal name + placeholders,
% and the SKELETON banner stays on every page.
\documentclass[10pt,twocolumn,a4paper]{article}
\usepackage[utf8]{inputenc}
\usepackage[T1]{fontenc}
\usepackage{lmodern}
\usepackage{amsmath,amssymb}
\usepackage{booktabs}
\usepackage[top=2.4cm,bottom=2.2cm,left=1.9cm,right=1.9cm]{geometry}
\setlength{\columnsep}{6mm}
\usepackage{xcolor}
\usepackage{microtype}            % protrusion + expansion: prevents overfull boxes
\setlength{\emergencystretch}{3em}
\usepackage{titlesec}
\titleformat{\section}{\bfseries}{\thesection}{0.5em}{}
\titleformat{\subsection}{\bfseries}{\thesubsection}{0.5em}{}
\titlespacing*{\section}{0pt}{1.5ex plus .3ex}{0.7ex}
\titlespacing*{\subsection}{0pt}{1.1ex plus .2ex}{0.5ex}
\titlespacing*{\paragraph}{0pt}{0.9ex plus .2ex}{0.5em}
\usepackage{enumitem}
\setlist{topsep=2pt,itemsep=2pt,parsep=0pt,leftmargin=1.3em}
\usepackage{hyperref}
\hypersetup{hidelinks,pdftitle={The Two Questions of Trust}}

% Skeleton banner so readers cannot mistake this for an accepted paper.
\usepackage{fancyhdr}
\pagestyle{fancy}\fancyhf{}
\fancyhead[C]{\footnotesize\color{gray}STRUCTURE SKELETON: not for circulation; experiments and results pending}
\fancyfoot[C]{\thepage}
\renewcommand{\headrulewidth}{0pt}

\definecolor{pendingclr}{HTML}{B00020}
\definecolor{writeclr}{HTML}{1A5FB4}
\newcommand{\priorfact}[1]{#1}
\newcommand{\pending}[1]{{\small\sffamily\color{pendingclr}[\,PENDING:\ #1\,]}}
\newcommand{\todo}[1]{{\small\sffamily\color{writeclr}[\,WRITE:\ #1\,]}}
\newcommand{\kwsep}{\ $\cdot$\ }

\begin{document}

\twocolumn[{%
\vspace*{-0.4cm}
{\footnotesize International Journal on Document Analysis and Recognition (IJDAR)\\
\ttfamily [DOI assigned on acceptance]\par}
\vspace{6pt}
\noindent\colorbox{gray!30}{\strut\ \footnotesize\bfseries ORIGINAL PAPER\ }%
\hfill{\footnotesize\color{gray}\itshape submission template}
\vspace{14pt}

{\noindent\LARGE\bfseries The Two Questions of Trust:\\[2pt]
Why a Model's Confidence Is Not Enough, and What Is\par}
\vspace{12pt}

{\noindent\large Authors Anonymized\par}
\vspace{4pt}
{\noindent\footnotesize\color{gray}STRUCTURE SKELETON. Submission draft only:
received / revised / accepted dates, DOI, and copyright line are assigned on
acceptance and are deliberately not shown.\par}
\vspace{16pt}

\begin{minipage}{\textwidth}\small
\noindent\textbf{Abstract}\quad
A deployed document-extraction model is trusted through a single number it
reports about itself: its own confidence. That number has a fatal confound,
it cannot separate an input that is hard but familiar from one that is easy
but drawn from a shifted world, and no amount of recalibration removes it,
because one scalar is being asked two different questions at once. In this
work we argue that trust in numeric document extraction is not one quantity
but two independent questions: is this answer internally consistent (a
symbolic question), and is the model still operating in the distribution it
was trained on (a statistical question). Two prior studies established,
separately, a strong answer to each question: a subset-sum arithmetic gate
that is precision-orthogonal to softmax confidence, and the variance of the
beam margin as a distribution-shift signal that confidence is structurally
blind to. Each is a complete result, and each is, by its authors' own
account, deliberately partial. The contribution of this paper is the
synthesis those papers could not make individually: a formal two-axis theory
of extraction trust, a non-redundancy result that explains mechanistically
why the two signals fail independently rather than together, and an
integrated, pre-registered evaluation testing whether composing them is
strictly safer than either axis or than confidence alone, at matched
false-alarm rate and matched human-review cost. We also state, unsoftened,
the boundary the synthesis does not move: the distributional axis remains a
population-level signal, not a per-receipt detector, and the framework is an
applied formalisation, not a methodological novelty. The paper's claim is
therefore sharp by construction: trust was never one number, it is two
questions, and asking both is the safer policy precisely where asking one is
not.
\par\vspace{9pt}
\noindent\textbf{Keywords}\quad selective prediction \kwsep distribution
shift \kwsep document total extraction \kwsep trustworthy machine learning
\kwsep structural verification
\end{minipage}
\vspace{1.5em}
}]

\paragraph{Contributions.}
\begin{enumerate}
  \item \textbf{A two-axis theory of extraction trust (new).} We formalise
  symbolic checkability and distributional stability as orthogonal trust
  axes, define their composition operator and the accept / abstain /
  flag-shift decision policy, and give a non-redundancy statement argued both
  definitionally and empirically.
  \item \textbf{A mechanistic account of why the axes are independent
  (new).} We show the symbolic axis's errors are governed by document
  structure while the distributional axis's signal is a scale (not location)
  effect, and tie this to an error-decorrelation analysis, turning an
  empirically convenient pairing into a predicted one.
  \item \textbf{An integrated system and a pre-registered head-to-head
  evaluation (new).} A shared, released benchmark and a pre-registered
  protocol testing the composed policy against each axis alone and against
  confidence alone, at fixed false-alarm and fixed reviewer cost, with
  operating-point analysis.
  \item \textbf{Hardened external validity for each axis (new evidence;
  component results are prior work).} New experiments extend the two
  established components; the base results for each axis are not re-claimed
  here, they are cited to prior work and used as the components this paper
  composes.
  \item \textbf{An explicit, carried-forward boundary (new framing).} We
  restate the permanent structural limitations of the components in the
  language of the framework, so the synthesis sharpens rather than overstates
  the prior claims.
\end{enumerate}

\section{Introduction: the crack}
\label{sec:intro}
\todo{Open on the confound, not the method: confidence cannot separate
hard-but-familiar from easy-but-shifted. State the two-question thesis.
Contributions paragraph. Then the explicit new-vs-prior delineation pointer
(see Appendix~\ref{app:delineation}) so this remains a legitimate extended
version that does not dilute the two prior papers.}

\section{Related work: why no single signal closes the crack}
\label{sec:related}
\todo{Selective prediction; confidence calibration; OOD / distribution-shift
detection; formal and structural verification; document KIE / total
extraction. The argument: each prior line maps to ONE axis only; none spans
both. This is what licenses the reframe.}

\section{The two-axis theory of trust}
\label{sec:theory}
\todo{Core, fully new. Definitions:}
\begin{itemize}
  \item Axis A, symbolic checkability: \todo{formal definition.}
  \item Axis B, distributional stability: \todo{formal definition.}
  \item Composition operator and the accept / abstain / flag-shift policy.
  \item Non-redundancy statement: argued definitionally here, tested
  empirically in Sec.~\ref{sec:mechanism} and Sec.~\ref{sec:integrated}.
\end{itemize}

\section{Axis A: symbolic verification, and its honest ceiling}
\label{sec:axisA}
\paragraph{Recap (prior work, not re-claimed).}
\priorfact{The subset-sum arithmetic gate is precision-orthogonal to softmax
confidence; on the pooled prior-work evaluation it reaches 0.989 precision
(Wilson lower bound 0.961), with either signal alone below 0.92.}
This result is established in prior work and cited; it is used here only as
the component this paper composes.
\paragraph{Hardening (new evidence).}
\begin{itemize}
  \item Alternative-verifier bake-off: \pending{precision/coverage/orthogonality vs. line-item, subtotal+tax, rounding verifiers}.
  \item Full precision--coverage--cost frontier: \pending{swept tolerance curves, per corpus and pooled, with CIs and the reviewer-cost model}.
  \item Power-resolving per-corpus replication: \pending{per-corpus CIs at pre-registered target half-width}.
  \item Real end-to-end latency on the deployment pipeline: \pending{true end-to-end ratio (prior work measured CPU-only standalone)}.
\end{itemize}

\section{Axis B: distributional stress, and its honest ceiling}
\label{sec:axisB}
\paragraph{Recap (prior work, not re-claimed).}
\priorfact{The variance of the beam margin is a distribution-shift signal
that is location-invariant where sequence confidence is not; on the prior
natural shift pair the variance-ratio statistic shows a large compression
that a Kolmogorov--Smirnov location test does not see.}
Established in prior work and cited; used here as a component only.
\paragraph{The boundary, carried forward unsoftened.}
\priorfact{The beam margin sits below every confidence baseline on raw
per-receipt AUROC and provides no per-receipt triage value. This is a
structural property, not a defect, and it is not retracted here.}
\paragraph{Hardening (new evidence).}
\begin{itemize}
  \item Multiple natural shift pairs: \pending{variance-ratio + significance battery across several natural pairs; how often sign/magnitude holds}.
  \item Controlled mechanism study: \pending{synthetic shifts with independently dialed difficulty vs. distribution distance}.
  \item Ablations and negative controls: \pending{beam width, length normalisation, architecture; registered no-shift null}.
  \item Operational monitor: \pending{detection rate at fixed false-alarm and detection delay (prior-work F6/F7 numbers were pending)}.
\end{itemize}

\section{The integrated system: the grand experiment}
\label{sec:integrated}
\todo{Shared, released benchmark and pipeline running both axes on the same
receipts.}
\begin{itemize}
  \item Pre-registered four-way head-to-head: composed vs. Axis A vs. Axis B
  vs. confidence alone, at matched false-alarm and matched reviewer cost.
  Primary endpoint: \pending{registered primary metric and result}.
  \item Blind-spot coverage test: on a CORD-only proxy (n=100, single
  beam-margin-paper pipeline) the composed policy (arith-pass AND
  $c_{\mathrm{seq}}\!\ge\!$ median) rejects 10 receipts that the symbolic
  axis alone wrongly accepts and 2 that confidence alone wrongly accepts;
  composed precision 0.83 at 0.36 coverage vs.\ confidence-alone 0.84 at
  0.50 and Axis-A-alone 0.70 at 0.54 (CORD-only proxy, n=100; exact-match
  on stored normalised fields, not the official KIE scorer). The integrated
  multi-corpus four-way version remains \pending{registered primary metric
  and result}.
  \item Error-decorrelation measurement: on the same CORD-only proxy
  (applicable receipts n=83) the Axis-A error event and the low-confidence
  event show $\phi=+0.29$ (permutation $p=0.011$, two-sided): a weak but
  significant positive association, not full independence, on this single
  pipeline (CORD-only proxy, n=83). This does not displace the cited
  prior-work orthogonality result, which rests on its own multi-corpus
  evidence; the integrated cross-corpus decorrelation analysis is
  \pending{empirical correlation of the two axes' error events at benchmark scale}.
  \item Adversarial / stress probe: \pending{graceful degradation of the composed policy under money-line noise + cardinality ablation}.
\end{itemize}

\section{Mechanism: why the coverage is principled, not lucky}
\label{sec:mechanism}
\todo{Scale-vs-location separation; why symbolic errors decorrelate from
confidence errors; the formal plus empirical account that turns a convenient
pairing into a predicted one. This section is what makes the story robust
rather than anecdotal, and it carries the novelty weight.}

\section{Limitations: the boundary, stated as strength}
\label{sec:limitations}
\begin{itemize}
  \item \priorfact{Axis B is a population-level signal, not a per-receipt
  detector; no per-receipt triage. Carried forward, not erased.}
  \item \priorfact{The framework is an applied formalisation; no
  methodological-novelty claim.}
  \item Residual external-validity bounds after the new evidence:
  \pending{state precisely what the broadened experiments do and do not establish}.
\end{itemize}

\section{Conclusion}
\label{sec:conclusion}
\todo{Restate the reframe: trust was never one number, it is two questions,
and you must ask both.}

\appendix
\section{New-vs-prior delineation}
\label{app:delineation}
\todo{Reproduce the table from DELINEATION.md so the extended-version
boundary is explicit in the manuscript itself.}

\section{Pre-registration record}
\label{app:prereg}
\todo{Reference PREREGISTRATION.md: hypotheses, endpoints, analysis plan,
fixed before any new experiment is run.}

\bibliographystyle{plain}
\bibliography{references}

\end{document}
